% !TeX program = lualatex
% halo:begin
% {
%   "version": 1,
%   "publish": true,
%   "course": "mathematical-analysis-i",
%   "section": "1.1",
%   "title": "数学分析（一）1.1：实数",
%   "categories": [
%     "study-notes"
%   ],
%   "tags": [
%     "mathematics"
%   ],
%   "excerpt": "习题 1.1 的课后作业：无理数、绝对值不等式与实数大小比较。"
% }
% halo:end
% 编译：lualatex 01-01-real-numbers.tex（建议运行两次）
% 课后作业根据《笔记 2027.pdf》第 1--3 页及习题 1.1 截图整理。
% LuaLaTeX 嵌入中文字体及 Unicode 映射，兼容 PDF 预览与文字复制。
\RequirePackage{iftex}
\RequireLuaTeX
\documentclass[UTF8, fontset=fandol, a4paper, 12pt]{ctexart}

\usepackage[margin=2.5cm]{geometry}
\usepackage{amsmath, amssymb, amsthm}
\usepackage{enumitem}
\usepackage{fancyhdr}
\usepackage[hidelinks]{hyperref}

\newcommand{\R}{\mathbb{R}}
\newcommand{\Nplus}{\mathbb{N}_{+}}
\theoremstyle{definition}
\newtheorem{definition}{定义}[section]
\newtheorem{proposition}[definition]{命题}
\newtheorem{theorem}[definition]{定理}
\newtheorem{example}[definition]{例题}
\renewcommand{\proofname}{证明}
\setlist[enumerate]{leftmargin=2em, itemsep=0.2em, topsep=0.3em}
\setlength{\headheight}{16pt}
\pagestyle{fancy}
\fancyhf{}
\fancyhead[L]{\small 数学分析（一）}
\fancyhead[R]{\small 1.1\quad 实数}
\fancyfoot[C]{\thepage}
\renewcommand{\headrulewidth}{0.3pt}
\raggedbottom

\title{实数}
\author{Yuchen Zhang}
\date{整理日期：2026年9月15日}
\makeatletter
\renewcommand{\maketitle}{%
  \begin{center}
    {\LARGE\bfseries \@title\par}\vspace{0.6em}
    \@author\par\vspace{0.2em}
    {\small \@date\par}
  \end{center}
}
\makeatother

\begin{document}
\maketitle

\section{课后作业}

\subsection{习题 1.1，第 1 题}
\noindent\textbf{题目：}设 $a\in\mathbb{Q}$，$x\in\R\setminus\mathbb{Q}$。证明：
（1）$a+x$ 为无理数；（2）当 $a\ne0$ 时，$ax$ 为无理数。

\begin{proof}
（1）反设 $a+x\in\mathbb{Q}$。由有理数对减法封闭，
$x=(a+x)-a\in\mathbb{Q}$，与 $x$ 为无理数矛盾。

（2）反设 $ax\in\mathbb{Q}$。由于 $a\in\mathbb{Q}$ 且 $a\ne0$，
有 $x=(ax)/a\in\mathbb{Q}$，仍与条件矛盾。
\end{proof}

\subsection{习题 1.1，第 3 题}
\noindent\textbf{题目：}设 $a,b\in\R$。若对任意 $\varepsilon>0$，都有 $|a-b|<\varepsilon$，证明 $a=b$。

\begin{proof}
反设 $a\ne b$，则 $|a-b|>0$。取 $\varepsilon=|a-b|$，由题设得到
$|a-b|<|a-b|$，矛盾。因此 $a=b$。
\end{proof}

\subsection{习题 1.1，第 5 题（1）}
\noindent\textbf{题目：}证明：对任意 $x\in\R$，有 $|x-1|+|x-2|\ge1$，并说明等号成立的条件。

\begin{proof}
由三角不等式，
\[
 |x-1|+|x-2|=|x-1|+|2-x|
 \ge |(x-1)+(2-x)|=1.
\]
等号成立当且仅当 $(x-1)(2-x)\ge0$，即 $x\in[1,2]$。
\end{proof}

\newpage
\subsection{习题 1.1，第 7 题}
\noindent\textbf{题目：}设 $x>0$，$b>0$，$a\ne b$。证明 $\dfrac{a+x}{b+x}$ 严格介于 $1$ 与 $\dfrac ab$ 之间。

\begin{proof}
由于 $b>0$、$b+x>0$，有
\[
 \frac{a+x}{b+x}-1=\frac{a-b}{b+x},\qquad
 \frac ab-\frac{a+x}{b+x}=\frac{x(a-b)}{b(b+x)}.
\]
当 $a>b$ 时，两式右端均为正，故
\[
 1<\frac{a+x}{b+x}<\frac ab.
\]
当 $a<b$ 时，两式右端均为负，故
\[
 \frac ab<\frac{a+x}{b+x}<1.
\]
\end{proof}

\subsection{习题 1.1，第 8 题}
\noindent\textbf{题目：}设 $p$ 为正整数且不是完全平方数，证明 $\sqrt p$ 是无理数。

\begin{proof}
反设 $\sqrt p=m/n$，其中 $m,n\in\Nplus$ 且 $\gcd(m,n)=1$。平方得
\[
 m^2=pn^2.
\]
若 $n>1$，取 $n$ 的一个素因子 $q$。由 $q\mid n$ 得 $q\mid m^2$；
由于 $q$ 是素数，可得 $q\mid m$。这与 $\gcd(m,n)=1$ 矛盾。
因此 $n=1$，从而 $p=m^2$，与 $p$ 不是完全平方数矛盾。
故 $\sqrt p$ 是无理数。
\end{proof}

\end{document}
